\section{La place des B-splines, des mélanges gaussiens et des réseaux de neurones}\label{sec-alternatives}

\subsection{Ce que les coefficients B-spline ne voient pas}
\begin{proposition}[Levier d'un pic]\label{prop-leverage}
Soit $B\in\R^{K\times q}$ la matrice de conception B-spline aux piliers et $H=B(B^\top B)^{-1}B^\top$ la matrice chapeau.
L'ajustement par moindres carrés d'un pic unitaire $e_k$ a une énergie $\|He_k\|^2=h_{kk}$, et $\sum_kh_{kk}=q$. Une
fraction $1-h_{kk}$ de l'énergie du pic reste dans le résidu et est invisible pour toute variable calculée à partir des
coefficients.
\end{proposition}
\begin{proof}
$H$ est symétrique et idempotente, donc $\|He_k\|^2=e_k^\top H^\top He_k=e_k^\top He_k=h_{kk}$, et
$\sum_kh_{kk}=\operatorname{tr}H=\rank B=q$. Les coefficients $(B^\top B)^{-1}B^\top e_k$ déterminent $He_k$ et rien d'autre.
\end{proof}

\begin{example}
Pour les six B-splines cubiques (nœuds en $\ln T$) du banc d'essai et les onze piliers,
\begin{center}\input{generated/\lang/tab_leverage}\end{center}
Aux piliers intérieurs, entre neuf mois et quatre ans, les coefficients ne captent que \vLevMin{} à \vLevMax{} de
l'énergie d'un pic (en moyenne $6/11\approx\num{0.545}$ sur l'ensemble des piliers). La spline répartit en outre ce qu'elle
capte sur plusieurs coefficients. Une anomalie sur un seul pilier, le motif de courbe aberrante le plus fréquent, est donc
diluée avant même que le regroupement ne commence.
\end{example}

\subsection{Pourquoi un mélange sur l'écart signale les erreurs de marquage FO}
\begin{proposition}\label{prop-gmm-fo}
Soient les variables les coefficients spline $c=B^+\Delta b$ de l'écart de repo. Un décalage parallèle du repo FO de
$\delta$ les déplace de $-\delta\,B^+\one$, une direction fixe. Si le mélange est ajusté sur un historique sain, où FO et Totem
concordent au bruit près, sa densité décroît dans cette direction : son taux d'alerte sur les erreurs de marquage FO tend
vers un quand $|\delta|$ augmente, quelle que soit la qualité de la courbe Totem.
\end{proposition}
\begin{proof}
$c$ est linéaire en $\Delta b$, et $\Delta b$ se déplace de $-\delta\one$. Chaque composante gaussienne a une variance finie dans
la direction $B^+\one$, donc la densité du mélange, et avec elle le score de vraisemblance, tend vers zéro dans cette
direction.
\end{proof}
Le détecteur structurel n'a pas ce problème : $T\cdot\one$ appartient au sous-espace du modèle nul, donc un écart de niveau
est retiré de chaque test et publié via $\hat\beta_0$. Le banc d'essai confirme les deux énoncés : \vFoGmmDiff{} des
erreurs de marquage FO sont signalées par le mélange sur l'écart, \vFoStruct{} par le détecteur structurel.

\subsection{Là où ils sont utiles}
\begin{itemize}
\item \textbf{P-splines dans le modèle nul.} Si $\beta_0+\beta_1\ln T$ se révèle trop rigide sur données réelles, une
spline pénalisée~\cite{eilers1996} peut le remplacer comme partie lisse du modèle nul ; les tests de la
\cref{sec-detection} restent inchangés.
\item \textbf{Mélanges sur des sorties structurées.} Un mélange gaussien sur les sorties du détecteur ($\hat\beta_0$, $\hat\beta_1$,
les $z_t$, $Q$ et les rapports de dispersion) apprend les régimes habituels par région et par catégorie
dans un espace de faible dimension où chaque coordonnée a un sens.
\item \textbf{Région $\times$ catégorie} sert au mieux à regrouper la calibration de $\kappa$ et de l'échelle de niveau,
pas à former des groupes.
\item \textbf{Réseaux de neurones.} D'après la \cref{prop-nn}, sur les mécanismes que l'on sait nommer, un réseau ne peut
au mieux qu'apprendre le filtre adapté. Les réseaux trouvent leur place une fois qu'existe un historique d'incidents
étiquetés et des motifs que les gabarits ne couvrent pas ; ils doivent alors être entraînés sur les résidus du modèle
structurel et sur des injections synthétiques, pas sur les courbes brutes.
\end{itemize}
